%-------------------------
% Resume Template in Latex
% Author : Renxu Bao 
% Style : TikZ Header + Dot Bullets
% License : MIT
%------------------------

\documentclass[a4paper,11pt]{article}

% 1. 中文支持 (必须使用 XeLaTeX 编译)
\usepackage[UTF8]{ctex}

% 2. 绘图包 (实现蓝色梯形标题)
\usepackage{tikz}
\usetikzlibrary{shapes, backgrounds}

% 3. 布局与格式包
\usepackage{latexsym}
\usepackage[empty]{fullpage}
\usepackage[explicit]{titlesec}
\usepackage{marvosym}
\usepackage{verbatim}
\usepackage{enumitem}
\usepackage[hidelinks]{hyperref}
\usepackage{fancyhdr}
\usepackage{tabularx}
\usepackage[usenames,dvipsnames]{xcolor}

%-------------------------------------------
% 颜色定义
\definecolor{MyDarkBlue}{RGB}{0, 50, 100}

% 页边距设置 (优化后让页面更饱满)
\pagestyle{fancy}
\fancyhf{} 
\fancyfoot{}
\renewcommand{\headrulewidth}{0pt}
\renewcommand{\footrulewidth}{0pt}

\addtolength{\oddsidemargin}{-0.6in}
\addtolength{\evensidemargin}{-0.6in}
\addtolength{\textwidth}{1.2in}
\addtolength{\topmargin}{-.7in}
\addtolength{\textheight}{1.4in}

\urlstyle{same}
\raggedbottom
\raggedright
\setlength{\tabcolsep}{0in}

% --- 自定义 TikZ 标题样式 ---
\titleformat{\section}
  {} 
  {} 
  {0pt} 
  {%
    \noindent 
    \begin{tikzpicture}[baseline=(t.base)]
    \vspace{-2pt}
      \node[inner xsep=10pt, inner ysep=3pt, text=white, font=\large\bfseries, align=left, anchor=base west] (t) at (0,0) {#1};
      \begin{scope}[on background layer]
         \fill[MyDarkBlue] (t.north west) -- (t.north east) -- ([xshift=8pt]t.south east) -- (t.south west) -- cycle;
         \draw[MyDarkBlue, line width=2pt] ([xshift=-68.30pt]t.south east) -- (\linewidth, 0 |- t.south east);
      \end{scope}
    \end{tikzpicture}%
  }
  []

\titlespacing*{\section}{0pt}{12pt}{6pt}

%-------------------------------------------
% 自定义命令

\newcommand{\resumeItem}[1]{
  \item\small{{#1 \vspace{-1pt}}}
}

\newcommand{\resumeEducationHeading}[4]{
  \vspace{-2pt}\item
    \begin{tabular*}{0.97\textwidth}[t]{l@{\extracolsep{\fill}}r}
      \textbf{#1} & #2 \\
      \textit{\small #3} & \textit{\small #4} \\
    \end{tabular*}\vspace{-5pt}
}

\newcommand{\resumeResearchHeading}[2]{
  \vspace{-2pt}\item
    \begin{tabular*}{0.97\textwidth}[t]{l@{\extracolsep{\fill}}r}
      \textbf{#1} & \textbf{#2} \\
    \end{tabular*}\vspace{-7pt}
}

\newcommand{\resumeProjectHeading}[2]{
    \item
    \begin{tabular*}{0.97\textwidth}{l@{\extracolsep{\fill}}r}
      \small#1 & #2 \\
    \end{tabular*}\vspace{-7pt}
}

\newcommand{\resumeSubHeadingListStart}{\begin{itemize}[leftmargin=0.15in, label={}]}
\newcommand{\resumeSubHeadingListEnd}{\end{itemize}}
\newcommand{\resumeItemListStart}{\begin{itemize}[leftmargin=*, label={\small\textbullet}]}
\newcommand{\resumeItemListEnd}{\end{itemize}\vspace{-2pt}}

%-------------------------------------------
%%%%%%  简历正文  %%%%%%%%%%%%%%%%%%%%%%%%%%%%

\begin{document}

%---------- 个人信息 ----------
\begin{center}
     \textbf{\Huge \scshape {\ziju{0.2} 你的姓名}} \\ \vspace{8pt}
    \small 
    \href{mailto:yourname@example.edu.cn}{\underline{yourname@example.edu.cn}} $|$ 
    +86 123-4567-8910 $|$ 
    \href{https://github.com/yourusername}{\underline{github.com/yourusername}} $|$
    \href{https://yourhomepage.com}{\underline{个人主页/Portfolio}}
\end{center}

%----------- 教育背景 -----------
\section{教育背景}
  \resumeSubHeadingListStart
    \resumeEducationHeading
      {厦门大学 (985 / 双一流)}{202X.09 -- 202X.06}
      {专业：计算机科学与技术}{厦门，福建}
      \resumeItemListStart
        \resumeItem{\textbf{GPA: 3.XX / 4.0} \ \ $|\ \ $ \textbf{专业排名: X / XXX (Top X\%)}}
        \resumeItem{\textbf{英语水平：}雅思 X.X / CET-6 XXX}
        \resumeItem{\textbf{研究兴趣：}人工智能、计算机视觉、多模态学习}
      \resumeItemListEnd
  \resumeSubHeadingListEnd

%----------- 技能清单 (有效增加页面垂直占据度) -----------
\section{个人技能}
 \begin{itemize}[leftmargin=0.15in, label={}]
    \small{\item{
     \textbf{专业技能}{: 深度学习框架 (PyTorch/TensorFlow), 计算机视觉, 跨模态表征学习} \\
     \textbf{编程/语言}{: Python (熟练), C/C++, SQL, LaTeX (精通), 英语 (雅思 7.0/CET-6)} \\
     \textbf{工程工具}{: Linux 环境开发, Docker 容器化, Git 版本控制, AWS/GCP 云计算平台}
    }}
 \end{itemize}

%----------- 科研经历 -----------
\section{科研经历}
  \resumeSubHeadingListStart

    \resumeResearchHeading
      {厦门大学媒体分析与计算组 (MAC)} 
      {202X.08 -- 202X.01}
      \resumeItemListStart
        \resumeItem{\textbf{项目：基于 AIGC 的跨模态图像增强研究} \hfill \textbf{\color{MyDarkBlue}[第一作者]}}
        \resumeItem{针对XX复杂环境下图像特征丢失问题，提出了一种基于稳定扩散 (Stable Diffusion) 的增强框架。}
        \resumeItem{设计了全新的对比学习损失函数，在保证语义一致性的同时，将 PSNR 指标提升了 15\%。}
        \resumeItem{该成果已投稿至 \textbf{CVPR 202X} 并获得审稿人正面评价。}
      \resumeItemListEnd

    \resumeResearchHeading
      {清华大学智能产业研究院 (AIR)} 
      {202X.02 -- 202X.07}
      \resumeItemListStart
        \resumeItem{\textbf{项目：面向智慧城市的自动驾驶场景建模} \hfill \textbf{\color{MyDarkBlue}[研究实习生]}}
        \resumeItem{负责处理大规模雷达点云数据与视觉信息的多模态对齐，优化了感知算法在雨天环境下的鲁棒性。}
        \resumeItem{主笔撰写了项目的技术报告，并参与了国家级重大科研课题的申报工作。}
      \resumeItemListEnd

    \resumeResearchHeading
      {中国政法大学 (跨学科交叉研究项目)} 
      {202X.09 -- 202X.12}
      \resumeItemListStart
        \resumeItem{\textbf{项目：法律文本的自动语义分析与知识图谱构建} \hfill \textbf{\color{MyDarkBlue}[核心成员]}}
        \resumeItem{利用大语言模型 (LLM) 对非结构化法律判决书进行实体抽取，构建了包含 10 万+节点的知识图谱。}
        \resumeItem{实现了跨学科背景下的技术落地，为法律人工智能助手提供了核心语义支持。}
      \resumeItemListEnd

  \resumeSubHeadingListEnd

%----------- 项目经历 -----------
\section{项目经历}
  \resumeSubHeadingListStart
    \resumeResearchHeading
      {项目/竞赛名称} 
      {202X.03 -- 202X.06}
      \resumeItemListStart
        \resumeItem{\textbf{获得奖项：某某大赛一等奖 (Top 1\%)} \hfill \textbf{\color{MyDarkBlue}[队长/核心开发]}}
        \resumeItem{项目背景描述：该系统旨在解决XX问题，填补了XX领域的空白。}
        \resumeItem{个人职责：负责系统架构设计、核心算法实现以及前后端联调。}
        \resumeItem{技术栈：UniApp, Python, PyTorch, Docker 等。}
      \resumeItemListEnd
  \resumeSubHeadingListEnd

%----------- 荣誉奖项 -----------
\section{荣誉奖项}
  \resumeSubHeadingListStart
    \resumeProjectHeading
      {\textbf{国家奖学金} (教育部颁发，Top 0.2\% 顶尖荣誉)}{202X}
    \resumeProjectHeading
      {\textbf{达美乐奖学金} (校级奖学金)}{202X}
    \resumeProjectHeading
      {\textbf{美国大学生数学建模竞赛 (MCM/ICM)} (Meritorious Winner)}{202X}
    \resumeProjectHeading
      {\textbf{中国高校计算机大赛 - AI 创意赛} (全国二等奖)}{202X}
  \resumeSubHeadingListEnd

\end{document}